\SourceLeaf{248}
\Article{373}
\textsc{Problem.} \emph{If the function $X$ has been resolved modulo $p$ into coprime factors $\xi,\xi',\xi''$, etc., resolve $X$ modulo $pp$ into similar factors $\Xi,\Xi',\Xi''$, etc., so that}
\[
\xi\equiv\Xi,\quad\xi'\equiv\Xi',\quad\xi''\equiv\Xi'',\text{ etc.}\;(\mathrm{mod.}\,p)
\]

\emph{Solution.} Let $X\equiv\xi\psi(\mathrm{mod.}\,p)$, or equivalently $X=\xi\psi+p\Sigma$. Put
\[
\Xi=\xi+p\varphi,\quad\Psi=\psi+p\omega
\]
Then
\[
\Xi\Psi=X-p\Sigma+(\varphi\psi+\xi\omega)p+pp\varphi\omega
\]
Thus, if $\Xi\Psi$ is to be $\equiv X(\mathrm{mod.}\,pp)$, $\varphi\psi+\xi\omega-\Sigma$ must necessarily be divisible by $p$. But since $\psi$ and $\xi$ are coprime functions modulo $p$, $\varphi$ and $\omega$ can be determined so that this condition is fulfilled (\textsection.\ 336), and moreover so that the degrees of $\varphi$ and $\omega$ are respectively one less than the degrees of the functions $\xi,\psi$. Hence $X\equiv\Xi\Psi(\mathrm{mod.}\,pp)$. It is clear that $\Psi$ can again be split in the same way into factors $\Xi'\Omega$, with the factor $\Xi'$ being $\equiv\xi'(\mathrm{mod.}\,p)$, and so on, until finally
\[
X\equiv\Xi\,\Xi'\,\Xi''\text{ etc.}\;(\mathrm{mod.}\,pp).\quad\text{Q. E. Fac.}
\]

\Needspace{20\baselineskip}\Article{374}
From this it is easy to prove that the function $X$ can also be resolved into factors modulo $p^3,p^4$, etc. More generally, let
\[
X\equiv PQ(\mathrm{mod.}\,p^m)\quad\text{or}\quad X=PQ+p^mR
\]
and let the function $P$ be coprime to $Q$ modulo $p$. On putting
\[
P'=P+Ap^m,\quad Q'=Q+Bp^m
\]
we have
\[
P'Q'=X-p^mR+(AQ+BP)p^m+ABp^{2m}
\]
Hence, for any modulus $p^\nu$ (with $\nu$ being $>m$ and $<2m+1$), we have
\[
P'Q'\equiv X,\quad\text{if}\quad R\equiv AQ+BP(\mathrm{mod.}\,p^{\nu-m})
\]

\SourceLeaf{249}
From this we see that, if the function $X$ has no equal divisors modulo $p$, it can be split into factors modulo $p^k$ in the same way as modulo $p$. But if $X$ has equal divisors, the matter becomes much more complicated and cannot even be fully dealt with by the preceding principles. Since we cannot communicate everything belonging to this subject, we shall consider only one case, which occurs very often and whose resolution is needed to settle certain doubts in what precedes. This is the case where only equal factors of degree one are considered. It can also be used directly to find the roots of congruences. We shall treat the matter in general on another occasion.

\Article{375}
Let, then, $X\equiv X'(x-a)^m(\mathrm{mod.}\,p)$, and let the function $X'$ be coprime to $x-a$: we seek all divisors of degree one congruent to $x-a$ modulo $p$ that divide $X$ modulo $pp,p^3$, etc. (We assume that the function $X$ is not divisible by $x-a$ without a modulus; otherwise $x-a$ would divide the function $X$ modulo every modulus.) If we substitute $z+a$ for $x$, we obtain
\[
Z\equiv Z'z^m(\mathrm{mod.}\,p)\quad\text{or}\quad Z=Z'z^m+pA
\]
Now, if $Z$ is divisible modulo $pp$ by some divisor of the form $z+\alpha p$, then $A$ must be of the form $zZ''+pB$. Unless this is so, the investigation is already finished. Let us therefore put
\[
\begin{aligned}
Z&\equiv Z'z^m+pZ''z(\mathrm{mod.}\,pp)\\
&\text{or}\quad Z=Z'z^m+pZ''z+ppB
\end{aligned}
\]
It is clear that $Z$ is divisible by $z$ and by every other divisor congruent to it modulo $p$.

To fix our attention, we put $m=4$; it is easy to see that every other case can be treated in a similar way. Now, if $Z$ is divisible modulo $p^3$ by some divisor of the form $z+\alpha p$, we have
\[
\begin{aligned}
0&\equiv-\alpha ppZ''+ppB(\mathrm{modd.}\,z+\alpha p,p^3)\\
&\text{or}\quad\alpha Z''\equiv B(\mathrm{modd.}\,z,p)
\end{aligned}
\]
There can now be three cases

1) If $Z''\equiv0(\mathrm{modd.}\,z,p)$ and $B$ is not $\equiv0$, then clearly no value of $\alpha$ satisfies the congruence, and hence $Z$ has, modulo $p^3$, no divisor of the form $z+\alpha p$. The investigation is therefore finished

\SourceLeaf{250}
\Needspace{14\baselineskip}2) If neither $Z''$ nor $B$ is $\equiv0(\mathrm{modd.}\,z,p)$, then $\alpha$ has a unique value, namely
\[
\alpha\equiv\frac B{Z''}(\mathrm{modd.}\,z,p).
\]
Thus there is a unique divisor $\equiv z+\alpha p(\mathrm{mod.}\,pp)$ of $Z$ modulo $p^3$, and we have
\[
Z\equiv V(z+\alpha p)+p^3W
\]
Now let the divisor of $Z(\mathrm{mod.}\,p^4)$ be $z+\alpha p+\beta pp$; then
\[
0\equiv\qquad\qquad
\]
\BookRule{.28\linewidth}
